\chapter{Related Works}

\section{Trajectory-Based Behavior Modeling in Traffic Research / Driving Style Identification}
\iffalse
Purpose: Show how maneuvers and trajectories are usually analyzed.

Discuss: Direct trajectory clustering, Driving style clustering, Cyclist behavior studies, Window-based feature extraction

Key contrast: Most approaches treat maneuvers as single homogeneous units or stop at window-level clustering.

Transition: Limited work performs hierarchical abstraction from local behavioral regimes to maneuver-level style representations.

traffic-studies: Orozco-Fontalvo
papers: kutsch, bachir , Belikhov, de zepeda, dias et al, milardo, Mobini Seraji, taubman
following: dias et al, kutsch, bachir, hoogendorm, zhang 2025 (cars), zhao 2017, 
overtaking, kutsch, Munnamgi, khan, mosacri (cars), mohammed 2019, rasch (cars), wang 2022, yin 2022

two approaches: car-following models, bicycle-following models 
closest

\fi

%%% section 1: riding style identification

\citeauthor{kutsch_analyzing_2025} (\citeyear{kutsch_analyzing_2025}) analyzed cyclist trajectories in urban intersections using the TUM-DOT MUC drone dataset. They identified following and overtaking maneuvers and characterized micro-level riding behavior using metrics such as speed, spacing, and headways. This work provides a high-resolution empirical foundation for trajectory- and maneuver-based analysis. Building on their dataset and maneuver definitions, the present study develops a modular, multi-stage clustering pipeline to identify latent riding styles across maneuvers, enabling hierarchical abstraction from local behavioral regimes to maneuver-level styles.

Several studies have focused on empirical analyses of cycling and driving behavior using trajectory or instrumented vehicle data. \citeauthor{belikhov_exploring_2025} (\citeyear{belikhov_exploring_2025}) analyzed free-flow cycling with instrumented bicycles in semi-controlled experiments, examining environmental and individual factors such as topography, wind, and trip characteristics. \citeauthor{dias_exploring_2023} (\citeyear{dias_exploring_2023}) investigated single-file cyclist trajectories, computing microscopic reaction delays and their dependence on spacing and relative speed. These studies provide detailed measurements of cyclist behavior in urban or controlled contexts and contribute quantitative insights into microscopic riding dynamics.

Other work has explored behavioral state and interaction modeling from trajectories. \citeauthor{mohammed_characterization_2019} (\citeyear{mohammed_characterization_2019}) applied multivariate finite mixture models to cyclist interactions, identifying following and overtaking states and quantifying transition patterns for microsimulation. \citeauthor{de_zepeda_dynamic_2021} (\citeyear{de_zepeda_dynamic_2021}) introduced dynamic clustering of vehicle trajectories to capture style switches over time. \citeauthor{zhang_shareable_2024} (\citeyear{zhang_shareable_2024}) proposed a hierarchical latent model to represent sequential driving behavior as fragments in low-dimensional space, extracting individual and shared driving styles in an unsupervised manner. Collectively, these studies demonstrate methods for segmenting, clustering, and representing behavioral states from trajectory data.

In addition to trajectory-based and behavioral state modeling, several studies have explored automated learning approaches for behavior recognition. Techniques such as clustering and classification frameworks have been applied to driving and cycling data to identify behavioral patterns and styles \cite{mobini_seraji_state---art_2025}. Other work has implemented AutoML pipelines for unsupervised feature extraction and behavior recognition, achieving high-accuracy classification on vehicle datasets \cite{milardo_data-driven_2022}. Deep learning approaches have also been combined with predictive modeling to link identified behavioral styles with maneuver-level outcomes, for example in lane-change risk assessment \cite{zhang_proactive_2023}. Collectively, these studies illustrate the growing application of automated and data-driven techniques for extracting patterns from large-scale trajectory data, complementing the more classical and interpretable approaches in cyclist and driving behavior research.

%%% section 2: overtaking / following

Several studies have examined following interactions among cyclists using trajectory data. \citeauthor{kutsch_analyzing_2025} (\citeyear{kutsch_analyzing_2025}) collected drone-based trajectories in urban environments and analyzed acceleration, speed, spacing, and headways for following and free-flow riding. Their work provides detailed empirical measurements of cyclist interactions, capturing heterogeneity over time, but remains descriptive without automated classification or predictive modeling of interactions. \citeauthor{dias_exploring_2023} (\citeyear{dias_exploring_2023}) conducted controlled experiments with single-file cyclists, quantifying reaction delays through correlation of acceleration and relative speed, highlighting the influence of spacing on following behavior. Similarly, \citeauthor{hoogendoorn_bicycle_2016} introduced a composite headway model to separate minimum and free headway distributions, capturing stochastic heterogeneity in following behavior. Complementary modeling efforts, such as \citeauthor{zhao_unified_2017} (\citeyear{zhao_unified_2017}) and \citeauthor{zhang_shareable_2024} (\citeyear{zhang_shareable_2024}), propose flow-based or AI-driven microscopic interaction models, providing a theoretical framework for follow-the-leader dynamics across traffic modes. Collectively, these studies establish descriptive and modeling baselines for following maneuvers, but primarily focus on descriptive statistics or parametric interaction models at a single analytical level, without hierarchical abstraction of local behavioral regimes into higher-level maneuver representations.

Overtaking behavior has been studied through both empirical measurements and algorithmic detection. \citeauthor{kutsch_analyzing_2025} (\citeyear{kutsch_analyzing_2025}) analyzed overtaking events in drone-collected trajectories, providing descriptive statistics on speed, lateral spacing, and maneuver duration. \citeauthor*{khan_characteristics_2001} (\citeyear{khan_characteristics_2001}) extended this approach with high-frequency tracking on constrained paths, quantifying passing dynamics and speed differences.
\citeauthor{mocsari_analysis_2009} (\citeyear{mocsari_analysis_2009}) collected qualitative observations on overtaking maneuvers to inform traffic microsimulation. 
Algorithmic approaches include \citeauthor{feng_automatic_2023} (\citeyear{feng_automatic_2023}), who proposed a multi-phase detection method for overtaking events using drone video data, and \citeauthor{wang_investigating_2022} (\citeyear{wang_investigating_2022}), who applied random parameter logit models to identify determinants of overtaking decisions.
While these studies provide descriptive, detection, or predictive capabilities, few integrate local behavioral states into maneuver-level style representations.



\section{Behavioral State Segmentation and Regime Discovery}
% motivation for states
Cycling behavior, like driving, unfolds continuously but can be meaningfully described through short-term, discrete behavioral units. Conceptual models of human action, such as those proposed by \citeauthor{pentland_modeling_1999} in \citetitle{pentland_modeling_1999}, suggest that sequences of micro-level states carry predictive information and reveal underlying structure. 

Trajectory segmentation for local behavioral state generation has been performed using event-based approaches, such as acceleration and braking \cite{yarlagadda_exploring_2022}, as well as sliding windows \cite{bachir_identifying_2025, li_volatility-based_2025, zhang_proactive_2023} and dynamic, context-aware windows determined by road alignment and maneuver content \cite{li_volatility-based_2025}.  Accordingly, the present work extracts volatility measures on riding dynamics within each window to define local riding regimes, ensuring that clusters reflect meaningful micro-behavioral variation rather than absolute magnitudes and providing a principled basis for subsequent aggregation and maneuver-level analysis.

% state characterization
% riding dynamics
In literature, microscopic driving dynamics are typically characterized using kinematic measures such as speed, acceleration (longitudinal and lateral), and yaw \cite{arvin_how_2019, arvin_safety_2021, bachir_identifying_2025, li_volatility-based_2025, wang_what_2015, yarlagadda_exploring_2022}. 
These measures capture fundamental aspects of motion and can be directly derived from trajectories. 
Longitudinal speed and acceleration provide insight into a rider's control along the direction of travel, while lateral acceleration and yaw reflect maneuvering and directional changes, which are crucial factors for understanding complex maneuvers and interactions.

% volatility
Volatility is a widely used metric for assessing driving risk, with studies in autonomous and human driving showing significant associations between driving volatility and crash frequencies across various contexts \cite{kamrani_extracting_2018, arvin_how_2019, zhang_proactive_2023}. 
Beyond crash prediction, volatility has also been effective in distinguishing and classifying different driving regimes and styles under diverse conditions \cite{mohammadnazar_classifying_2021, li_volatility-based_2025, yarlagadda_exploring_2022, wang_what_2015}.
However, scenario segmentation is critical for accurate volatility-based analysis, as driving volatilities and vehicle dynamics vary significantly across different contexts and traffic situations \cite{mohammadnazar_classifying_2021}, especially between intersection and non-intersection environments in urban settings \cite{li_volatility-based_2025}.

% traffic
While Volatility captures individual driving dynamics and behavioral risk, a broader set of surrogate safety measures is often deployed to capture the surrounding traffic and evaluate potential conflicts \cite{wang_review_2021}. These measures provide additional information about the operational and safety-relevant context of maneuvers, enabling a more comprehensive analysis of behavior under complex traffic conditions. Frequently used SSMs include Time-To-Collision (TTC), Deceleration-Rate-To-Avoid-Crash (DRAC), Distance-At-Closest-Approach (DCA) and proximity counts of other moving traffic participants.

Incorporating these surrogate safety measures allows relating behavioral regimes and maneuver-execution modes to the surrounding traffic context, supporting interpretability and an exploratory analysis.
%\cite{wang_review_2021, bahmankhah_interaction_2019, hu_high-resolution_2022, huang_predicting_2024, nabavi_niaki_is_2019}
% infrastructure

While traffic and infrastructure can influence behavior \cite{gindele_probabilistic_2010, sharif_context-awareness_2017, he_multi-scenario_2025}, the identification of local regimes in this work is performed independently of such situational factors. Contextual variables are instead analyzed post hoc to interpret observed behavioral patterns, providing descriptive insight into how maneuvers manifest under varying traffic and infrastructural conditions without affecting the unsupervised abstraction of micro-level states.

% state aggregation
Although probabilistic frameworks such as Markov Chains or Hidden Markov models have been used to model sequences of behavioral states \cite{khattak_analysis_2017, yin_extracting_2022, he_multi-scenario_2025} and personal driving styles \cite{meng_human_2006,zou_multivariate_2022}], the conceptual value lies in the recognition that sequences of discrete regimes encode structure: transition rates, dwell times, and entropy measures provide principled descriptors of behavioral complexity without requiring explicit generative assumptions \cite{vegetabile_estimating_2019}. In this sense, a maneuver can be understood as an ordered composition of behavioral regimes, whose statistical properties summarize internal variability and execution patterns. Post hoc analysis of environmental context further allows interpretation of how these behaviors manifest under different traffic or infrastructure conditions, enhancing insight without constraining the abstraction process itself.

Approaches that identify latent behavioral states provide a conceptual bridge from descriptive statistics to unsupervised style analysis. 
\citeauthor{mohammed_characterization_2019} (\citeyear{mohammed_characterization_2019}) applied Gaussian Finite Mixture Models to classify cyclist following and overtaking interactions into distinct motion states based on longitudinal distance, lateral distance, and speed differences. While this probabilistic framework enables flexible state identification, the analysis remains limited to maneuver-specific state classification and descriptive interpretation. The study thus exemplifies parametric state discovery at a single analytical level rather than multi-stage behavioral abstraction.
Extending this concept, \citeauthor{bachir_identifying_2025} (\citeyear{bachir_identifying_2025}) proposed a two-level unsupervised clustering framework to characterize cyclist riding styles under free-flow conditions. At the first level, time-window-based riding volatility measures derived from speed, acceleration, and rotational fluctuations were clustered using k-means to identify local behavioral regimes. At the second level, trajectory-level style representations were constructed by aggregating proportions of level-1 states, transition rates, entropy measures, and contextual features such as infrastructure exposure. These aggregated descriptors were subsequently reclustered to identify rider-level styles. While this approach demonstrates the potential of multi-level unsupervised clustering for behavioral characterization and incorporates stability considerations via silhouette and ARI evaluation, it is not maneuver-specific and does not explicitly frame the process as hierarchical abstraction within complex traffic interactions. Extending this concept toward maneuver-level regime aggregation and style discovery forms a central methodological motivation of the present study.

Taken together, this body of work motivates a modular, assumption-light abstraction of cyclist behavior: local behavioral regimes capture micro-level dynamics, their composition and distribution convey structural information about maneuvers, and contextual features provide descriptive interpretability. Conceptually, this approach aligns with prior research on driving style and behavior modeling while extending it to multi-stage, maneuver-level abstraction, providing a principled foundation for hierarchical behavioral analysis.



\section{Multi-Stage and Hierarchical Clustering Frameworks}
\iffalse
Discuss: Two-stage clustering, Hierarchical abstraction, Modular pipelines, Representation → clustering separation

Key argument: Existing approaches either perform hierarchical clustering within a single algorithmic framework or rely on end-to-end learning. Modular multi-stage clustering with explicit intermediate representation aggregation is less studied.

papers: nielsen, scheiner, ren 2017, lai, nakashima

\fi

Multi-stage clustering frameworks provide a structured approach to uncovering patterns in complex datasets by separating the analysis into modular steps. Rather than applying a single clustering algorithm directly to the raw data, these approaches first identify local or intermediate patterns, transform or aggregate them into higher-level representations, and then perform subsequent clustering to extract global structure. This strategy is particularly beneficial for data with heterogeneous, sequential, or high-dimensional characteristics, as it improves interpretability, enables stage-wise optimization, and allows flexible adaptation to different problem domains \cite{nielsen_hierarchical_2016, scheiner_multi-stage_2019}.

Applications of multi-stage clustering can be found across a variety of fields, including time series analysis, sensor data processing, and behavioral modeling. For example, approaches in temporal data segment time series into sliding windows, perform initial clustering within each segment, and then aggregate the results to derive global clusters \cite{ren_sliding_2017, nakashima_performance_2016, lai_novel_2010, guijo-rubio_time-series_2021}. Across these studies, the core principle remains consistent: first-stage clustering captures fine-grained, local patterns, while second-stage clustering abstracts these into higher-level structures, allowing both local variation and global patterns to be represented.

In the present study, this multi-stage paradigm is applied to cyclist trajectory data. First, short temporal windows are clustered to identify local riding regimes, which capture micro-level variations in speed, acceleration, and rotational behavior. These local regimes are then aggregated into maneuver-level descriptors, which summarize the internal structure of each maneuver. A second-stage clustering of these descriptors produces latent maneuver styles, providing hierarchical abstraction from local behavior to maneuver-level patterns. This design leverages the interpretability and modularity of multi-stage clustering while adapting it specifically to trajectory-based behavior analysis.


\section{Unsupervised Model Selection and Stability}
\iffalse
Discuss: Sensitivity of clustering to initialization, Silhouette limitations, Stability-based validation, 

Transition: These considerations motivate the stability-oriented model selection strategy adopted in this thesis.

papers: chen et al, coates and ng, 
\fi





Existing work either clusters full trajectories directly, identifies local behavioral regimes without higher-level abstraction, or relies on integrated end-to-end frameworks. A modular, interpretable, multi-stage clustering approach that first extracts local behavioral regimes and subsequently derives maneuver-level style representations remains insufficiently explored. This motivates the hierarchical framework proposed in the following chapter.
While these approaches demonstrate the effectiveness of unsupervised maneuver-level behavior analysis, existing methods either emphasize direct time-series clustering or do not explicitly separate local regime discovery from higher-level maneuver abstraction. A modular, stability-oriented multi-stage clustering framework that transforms local behavioral regimes into aggregated maneuver style representations remains insufficiently explored. This motivates the hierarchical pipeline introduced in the following chapter.